% Generated by make_book.py from Lecture3/appendix_topology.tex; edit the lecture, not this file.
\chapter{Topology and Differential Topology}\label{A-chapter}
% The group keeps settings such as \emergencystretch inside the chapter.
\begingroup
\section*{Introduction}
\addcontentsline{toc}{section}{Introduction}

In Lecture~3, we used several facts without proof: that a continuous function on a closed and bounded set attains its maximum, that a function with continuous partial derivatives is differentiable (Theorem~\ref{L3-thm:C1}), and the implicit function theorem in the proof of the Lagrange multiplier theorem (Theorem~\ref{L3-thm:lagrange}). We also met sets that are not open subsets of $\mathbb{R}^n$, but curved objects: spheres, the set of orthogonal matrices, and level sets of functions. This appendix provides the mathematical background.
\begin{itemize}
    \item Sections~\ref{A-sec:metric}--\ref{A-sec:connected} are about topology, the mathematics of nearness and continuity. Its central notions are open sets, compactness and connectedness. Topology studies the properties of sets that do not change under continuous deformations.
    \item Sections~\ref{A-sec:inverse}--\ref{A-sec:morse} are about differential topology, which studies manifolds, i.e., sets that look like $\mathbb{R}^k$ near each of their points, and differentiable maps between them. Its main tools are the inverse and implicit function theorems. At the end, Morse theory connects the critical points of a function, classified by the eigenvalues of its Hessian as in Lecture~3, with the shape of the space on which it is defined.
    \item Section~\ref{A-sec:ml} collects some places where these ideas appear in machine learning.
\end{itemize}
We use the notation of Lecture~3: vectors are bold, and $\norm{\cdot}$ is the Euclidean norm unless we say otherwise.

\section{Metric Spaces and Open Sets}\label{A-sec:metric}

\begin{definition}[Metric space]\label{A-def:metric}
    A metric on a set $X$ is a function $d: X \times X \to \mathbb{R}$ such that for all $x, y, z \in X$:
    \begin{enumerate}
        \item $d(x, y) \geq 0$, and $d(x, y) = 0$ if and only if $x = y$;
        \item $d(x, y) = d(y, x)$;
        \item $d(x, z) \leq d(x, y) + d(y, z)$ (triangle inequality).
    \end{enumerate}
    A set with a metric is a metric space.
\end{definition}

\begin{example}\label{A-ex:metrics}
    \begin{enumerate}
        \item Every norm on a vector space defines a metric $d(\mathbf{x}, \mathbf{y}) = \norm{\mathbf{x} - \mathbf{y}}$ (Lecture~2). In particular, $\mathbb{R}^n$ with the $\ell_1$, $\ell_2$ or $\ell_\infty$ norm is a metric space.
        \item On the space $C[a, b]$ of continuous functions on $[a, b]$, $d(f, g) = \max_{x \in [a, b]}|f(x) - g(x)|$ is a metric.
        \item Every subset of a metric space is a metric space with the same metric. For example, the unit sphere $S^{n-1} = \{\mathbf{x} \in \mathbb{R}^n : \norm{\mathbf{x}} = 1\}$ is a metric space with the Euclidean distance.
        \item On any set, the discrete metric $d(x, y) = 1$ for $x \neq y$ (and $d(x, x) = 0$) is a metric.
    \end{enumerate}
\end{example}

\begin{definition}[Open and closed sets]\label{A-def:open-closed}
    Let $(X, d)$ be a metric space. The open ball with center $x$ and radius $r > 0$ is $B(x, r) = \{y \in X : d(x, y) < r\}$. A set $U \subseteq X$ is open if for every $x \in U$ there is an $r > 0$ with $B(x, r) \subseteq U$. A set $C \subseteq X$ is closed if its complement $X \setminus C$ is open.
\end{definition}

For $X = \mathbb{R}^n$, this is the definition of open sets in Lecture~3. Closed is not the opposite of open: $\mathbb{R}^n$ and $\emptyset$ are both open and closed, and the half-open interval $[0, 1)$ is neither open nor closed in $\mathbb{R}$.

\begin{proposition}\label{A-prop:open-sets}
    In a metric space $X$:
    \begin{enumerate}
        \item every open ball is open;
        \item every union of open sets is open, and every intersection of finitely many open sets is open;
        \item every intersection of closed sets is closed, and every union of finitely many closed sets is closed.
    \end{enumerate}
\end{proposition}

\begin{proof}
    (1) If $y \in B(x, r)$, let $s = r - d(x, y) > 0$. For $z \in B(y, s)$, the triangle inequality gives $d(x, z) \leq d(x, y) + d(y, z) < d(x, y) + s = r$, so $B(y, s) \subseteq B(x, r)$.

    (2) If $x$ lies in a union of open sets, it lies in one of them, $U$, and some ball $B(x, r) \subseteq U$ lies in the union. If $x \in U_1 \cap \cdots \cap U_m$ with $B(x, r_i) \subseteq U_i$, then $B(x, \min_i r_i)$ lies in the intersection.

    (3) Take complements in (2).
\end{proof}

Infinite intersections of open sets need not be open: $\bigcap_{k \geq 1}(-\frac{1}{k}, \frac{1}{k}) = \{0\}$. For a set $A \subseteq X$, the interior $\interior A$ is the union of all open sets contained in $A$ (the largest open subset of $A$), the closure $\overline{A}$ is the intersection of all closed sets containing $A$ (the smallest closed set containing $A$), and the boundary is $\partial A = \overline{A} \setminus \interior A$. For example, the closed ball $\{\mathbf{x} : \norm{\mathbf{x}} \leq 1\}$ in $\mathbb{R}^n$ has the open ball as interior and the sphere $S^{n-1}$ as boundary.

\begin{definition}[Convergence]\label{A-def:convergence}
    A sequence $(x_k)$ in a metric space converges to $x$ if $d(x_k, x) \to 0$ as $k \to \infty$. It is a Cauchy sequence if for every $\varepsilon > 0$ there is an $N$ with $d(x_k, x_l) < \varepsilon$ for all $k, l \geq N$.
\end{definition}

Closed sets are exactly the sets that cannot be left by taking limits.

\begin{proposition}\label{A-prop:closed-limits}
    A subset $C$ of a metric space $X$ is closed if and only if the limit of every convergent sequence of points of $C$ lies in $C$.
\end{proposition}

\begin{proof}
    Let $C$ be closed, and let $x_k \in C$ converge to $x$. If $x \notin C$, then some ball $B(x, r)$ lies in the open set $X \setminus C$; but $x_k \in B(x, r)$ for large $k$, a contradiction. Conversely, suppose that $C$ is not closed. Then $X \setminus C$ is not open, so there is an $x \notin C$ such that every ball $B(x, \frac{1}{k})$ contains a point $x_k \in C$. Then $x_k \to x$, and $x \notin C$.
\end{proof}

\section{Continuity and Topological Spaces}\label{A-sec:continuity}

Continuity can be described with open sets alone, without mentioning distances.

\begin{proposition}\label{A-prop:continuity-open}
    Let $f: X \to Y$ be a map between metric spaces. Then $f$ is continuous (at every point, in the $\varepsilon$-$\delta$ sense of Lecture~3) if and only if for every open set $V \subseteq Y$, the preimage $f^{-1}(V) = \{x \in X : f(x) \in V\}$ is open.
\end{proposition}

\begin{proof}
    Let $f$ be continuous, $V \subseteq Y$ open, and $x \in f^{-1}(V)$. There is an $\varepsilon > 0$ with $B(f(x), \varepsilon) \subseteq V$, and by continuity a $\delta > 0$ such that $d(x, x') < \delta$ implies $d(f(x), f(x')) < \varepsilon$, i.e., $B(x, \delta) \subseteq f^{-1}(V)$. Conversely, if preimages of open sets are open, then for $x \in X$ and $\varepsilon > 0$, the set $f^{-1}(B(f(x), \varepsilon))$ is open and contains $x$, so it contains a ball $B(x, \delta)$; this is the $\varepsilon$-$\delta$ condition.
\end{proof}

Taking complements, $f$ is continuous if and only if preimages of closed sets are closed.

\begin{corollary}\label{A-cor:level-sets-closed}
    Let $f: \mathbb{R}^n \to \mathbb{R}$ be continuous and $c \in \mathbb{R}$. Then the set $\{f < c\}$ is open, and the sublevel set $\{f \leq c\}$ and the level set $\{f = c\}$ are closed. The same holds for intersections of finitely many such sets.
\end{corollary}

\begin{proof}
    These are the preimages of the open set $(-\infty, c)$ and the closed sets $(-\infty, c]$ and $\{c\}$.
\end{proof}

So the feasible sets of the constrained problems in Lecture~3, such as $\{\mathbf{x} : g_1(\mathbf{x}) \leq 0, \ldots, g_m(\mathbf{x}) \leq 0\}$ for continuous $g_i$, are closed. Similarly, the set $GL(n, \mathbb{R})$ of invertible $n \times n$ matrices is open in $\mathbb{R}^{n \times n}$, as the preimage of the open set $\mathbb{R} \setminus \{0\}$ under the continuous function $\det$ (a polynomial in the entries).

Proposition~\ref{A-prop:continuity-open} suggests an abstraction: forget the metric, and only remember which sets are open.

\begin{definition}[Topological space]\label{A-def:topology}
    A topology on a set $X$ is a collection $\mathcal{T}$ of subsets of $X$, called open sets, such that $\emptyset, X \in \mathcal{T}$, every union of sets in $\mathcal{T}$ lies in $\mathcal{T}$, and every intersection of finitely many sets in $\mathcal{T}$ lies in $\mathcal{T}$. A map $f: X \to Y$ between topological spaces is continuous if preimages of open sets are open. A homeomorphism is a continuous bijection whose inverse is also continuous; if one exists, $X$ and $Y$ are homeomorphic.
\end{definition}

By Proposition~\ref{A-prop:open-sets}, the open sets of a metric space form a topology. Two metrics can give the same topology: we will see in Theorem~\ref{A-thm:norms-equivalent} that all norms on $\mathbb{R}^n$ do. Homeomorphic spaces are the same from the point of view of topology: a homeomorphism matches up their open sets, convergent sequences, continuous functions, and so on.

\begin{example}\label{A-ex:homeomorphisms}
    \begin{enumerate}
        \item The sigmoid function is a homeomorphism from $\mathbb{R}$ to $(0, 1)$, with the continuous inverse $p \mapsto \log\frac{p}{1 - p}$ (the logit). Similarly, $\tanh$ is a homeomorphism $\mathbb{R} \to (-1, 1)$. So bounded and unbounded sets can be homeomorphic: lengths and distances are not topological properties.
        \item The open ball $B(\mathbf{0}, 1) \subseteq \mathbb{R}^n$ is homeomorphic to $\mathbb{R}^n$, via $\mathbf{x} \mapsto \frac{\mathbf{x}}{1 - \norm{\mathbf{x}}}$, with the inverse $\mathbf{y} \mapsto \frac{\mathbf{y}}{1 + \norm{\mathbf{y}}}$.
        \item A square and a disk in $\mathbb{R}^2$ are homeomorphic (rescale each ray from the center), and so are a circle and the boundary of a square: corners are not a topological property. But a circle and an interval are not homeomorphic (Example~\ref{A-ex:connected}).
    \end{enumerate}
\end{example}

\section{Compactness}\label{A-sec:compact}

The real numbers have no gaps. This is the completeness axiom, which we take as one of the defining properties of $\mathbb{R}$: every nonempty set of real numbers that is bounded above has a least upper bound (supremum). It implies that every bounded monotone sequence of real numbers converges (to the supremum or the infimum of its values).

\begin{theorem}[Bolzano--Weierstrass]\label{A-thm:bolzano}
    Every bounded sequence in $\mathbb{R}^n$ has a convergent subsequence.
\end{theorem}

\begin{proof}
    First let $n = 1$. Call an index $k$ a peak if $x_k \geq x_l$ for all $l > k$. If there are infinitely many peaks $k_1 < k_2 < \cdots$, then $x_{k_1} \geq x_{k_2} \geq \cdots$ is a monotone subsequence. Otherwise, there is an index $k_1$ after which there are no peaks; then $k_1$ is not a peak, so $x_{k_2} > x_{k_1}$ for some $k_2 > k_1$; $k_2$ is not a peak either, and so on. This gives an increasing subsequence. In both cases, the monotone subsequence is bounded, so it converges. For $n > 1$, choose a subsequence whose first coordinates converge, then a subsequence of it whose second coordinates converge, and so on; after $n$ steps all coordinates converge.
\end{proof}

\begin{corollary}\label{A-cor:complete}
    $\mathbb{R}^n$ is complete: every Cauchy sequence in $\mathbb{R}^n$ converges.
\end{corollary}

\begin{proof}
    A Cauchy sequence $(\mathbf{x}_k)$ is bounded (all terms after some $N$ lie within distance $1$ of $\mathbf{x}_N$), so by Theorem~\ref{A-thm:bolzano} a subsequence converges to some $\mathbf{x}$. For $\varepsilon > 0$, choose $N$ such that $\norm{\mathbf{x}_k - \mathbf{x}_l} < \varepsilon$ for $k, l \geq N$, and an element $\mathbf{x}_l$ of the subsequence with $l \geq N$ and $\norm{\mathbf{x}_l - \mathbf{x}} < \varepsilon$. Then $\norm{\mathbf{x}_k - \mathbf{x}} < 2\varepsilon$ for all $k \geq N$.
\end{proof}

\begin{definition}[Compact set]\label{A-def:compact}
    A subset $K$ of a metric space is compact if every sequence in $K$ has a subsequence that converges to a point of $K$.
\end{definition}

(In general topological spaces, one defines compactness with open covers: $K$ is compact if every collection of open sets that covers $K$ contains finitely many sets that already cover $K$. For metric spaces, the two definitions agree; we do not prove this.)

\begin{theorem}[Heine--Borel]\label{A-thm:heine-borel}
    A subset of $\mathbb{R}^n$ is compact if and only if it is closed and bounded.
\end{theorem}

\begin{proof}
    Let $K$ be closed and bounded, and let $(\mathbf{x}_k)$ be a sequence in $K$. By Theorem~\ref{A-thm:bolzano}, a subsequence converges, and by Proposition~\ref{A-prop:closed-limits}, its limit lies in $K$. Conversely, if $K$ is unbounded, there are $\mathbf{x}_k \in K$ with $\norm{\mathbf{x}_k} \geq k$, and no subsequence converges, since convergent sequences are bounded. If $K$ is not closed, there is a sequence in $K$ that converges to a point outside $K$ (Proposition~\ref{A-prop:closed-limits}), and all its subsequences converge to the same point, so none converges to a point of $K$.
\end{proof}

For example, spheres, closed balls (for any norm), and the set $O(n)$ of orthogonal matrices are compact. $O(n)$ is closed, as the preimage of $\{I\}$ under the continuous map $Q \mapsto Q^\top Q$, and bounded, since $\norm{Q}_F^2 = \tr(Q^\top Q) = n$. The hyperbola $\{xy = 1\}$ is closed but unbounded, and the open ball is bounded but not closed; neither is compact.

\begin{theorem}[Extreme value theorem]\label{A-thm:extreme-value}
    Let $K$ be a nonempty compact set, and let $f: K \to \mathbb{R}$ be continuous. Then $f$ attains a minimum and a maximum on $K$.
\end{theorem}

\begin{proof}
    Let $m = \inf\{f(\mathbf{x}) : \mathbf{x} \in K\}$, where we allow $m = -\infty$. Choose $\mathbf{x}_k \in K$ with $f(\mathbf{x}_k) \to m$. A subsequence converges to some $\mathbf{x}^* \in K$, and by continuity, $f(\mathbf{x}^*) = m$. So $m$ is finite, and it is the minimum. For the maximum, apply this to $-f$.
\end{proof}

This justifies the claims in Lecture~3 that $f(x, y) = x + y$ and the Rayleigh quotient $\mathbf{x}^\top A\mathbf{x}$ attain their maxima on the unit sphere (Examples~\ref{L3-ex:lagrange-circle} and~\ref{L3-ex:lagrange-eigen}). On $\mathbb{R}^n$, which is not compact, a minimum need not exist, but the following condition is often enough.

\begin{proposition}[Coercive functions]\label{A-prop:coercive}
    Let $f: \mathbb{R}^n \to \mathbb{R}$ be continuous and coercive, i.e., $f(\mathbf{x}) \to \infty$ as $\norm{\mathbf{x}} \to \infty$. Then $f$ has a global minimum.
\end{proposition}

\begin{proof}
    Fix any $\mathbf{x}_0$. The sublevel set $S = \{f \leq f(\mathbf{x}_0)\}$ is closed (Corollary~\ref{A-cor:level-sets-closed}) and bounded (by coercivity), so it is compact and nonempty, and $f$ has a minimum $\mathbf{x}^*$ on $S$. Points outside $S$ have values larger than $f(\mathbf{x}_0) \geq f(\mathbf{x}^*)$, so $\mathbf{x}^*$ is a global minimum.
\end{proof}

\begin{example}\label{A-ex:existence-ml}
    A loss $L(\mathbf{w}) \geq 0$ plus a ridge penalty, $L(\mathbf{w}) + \lambda\norm{\mathbf{w}}^2$ with $\lambda > 0$, is coercive, so it has a global minimum whenever $L$ is continuous. Without regularization, a minimum need not exist: for logistic regression with the single training example $x = 1$, $y = 1$, the loss $\log(1 + e^{-w})$ decreases towards $0$ as $w \to \infty$, but never reaches it. The same happens for every linearly separable data set, where scaling up a separating weight vector always decreases the loss.
\end{example}

Compactness also explains why all norms on $\mathbb{R}^n$ lead to the same notions of convergence and continuity, as claimed in Lecture~3.

\begin{theorem}[Equivalence of norms]\label{A-thm:norms-equivalent}
    For every norm $N$ on $\mathbb{R}^n$, there are constants $0 < c \leq C$ with $c\norm{\mathbf{x}} \leq N(\mathbf{x}) \leq C\norm{\mathbf{x}}$ for all $\mathbf{x}$. Consequently, all norms on $\mathbb{R}^n$ define the same open sets.
\end{theorem}

\begin{proof}
    Writing $\mathbf{x} = \sum_i x_i\mathbf{e}_i$, the triangle inequality and the Cauchy--Schwarz inequality give
    \[N(\mathbf{x}) \leq \sum_i |x_i| N(\mathbf{e}_i) \leq \Big(\sum_i N(\mathbf{e}_i)^2\Big)^{1/2}\norm{\mathbf{x}} = C\norm{\mathbf{x}}.\]
    In particular, $|N(\mathbf{x}) - N(\mathbf{y})| \leq N(\mathbf{x} - \mathbf{y}) \leq C\norm{\mathbf{x} - \mathbf{y}}$, so $N$ is continuous. By Theorem~\ref{A-thm:extreme-value}, $N$ attains a minimum $c$ on the compact unit sphere $S^{n-1}$, and $c > 0$ since $N$ is positive on $S^{n-1}$. For $\mathbf{x} \neq \mathbf{0}$, $N(\mathbf{x}) = \norm{\mathbf{x}}\, N\big(\frac{\mathbf{x}}{\norm{\mathbf{x}}}\big) \geq c\norm{\mathbf{x}}$. Finally, the inequalities show that every $N$-ball around a point contains a Euclidean ball around it and vice versa, so the open sets are the same.
\end{proof}

\begin{remark}[Infinite dimensions]
    In infinite-dimensional spaces, closed and bounded sets need not be compact, and different norms can give different topologies. For example, in the space of square-summable sequences with the $\ell_2$ norm, the unit vectors $\mathbf{e}_1, \mathbf{e}_2, \ldots$ lie in the closed unit ball, but $\norm{\mathbf{e}_j - \mathbf{e}_k} = \sqrt{2}$ for $j \neq k$, so no subsequence converges. This is one reason why analysis in function spaces, such as the spaces of functions in Lecture~2, is harder than in $\mathbb{R}^n$.
\end{remark}

\section{Connectedness}\label{A-sec:connected}

\begin{definition}[Connected sets]\label{A-def:connected}
    A subset $X$ of a metric space is disconnected if $X \subseteq U \cup V$ for two open sets $U$ and $V$ such that $X \cap U$ and $X \cap V$ are nonempty and disjoint. Otherwise, $X$ is connected. It is path-connected if any two points $x, y \in X$ can be joined by a path, i.e., a continuous map $\gamma: [0, 1] \to X$ with $\gamma(0) = x$ and $\gamma(1) = y$.
\end{definition}

\begin{proposition}\label{A-prop:connected}
    \begin{enumerate}
        \item Every interval of $\mathbb{R}$ is connected.
        \item If $X$ is connected and $f$ is continuous on $X$, then $f(X)$ is connected.
        \item Every path-connected set is connected.
    \end{enumerate}
\end{proposition}

\begin{proof}
    (1) Suppose an interval $I$ is covered by open sets $U$ and $V$ as in the definition, with $a \in I \cap U$ and $b \in I \cap V$; by symmetry, we may assume $a < b$. Then $[a, b] \subseteq I$. Let $s = \sup\{t \in [a, b] : [a, t] \subseteq U\}$; then $[a, t] \subseteq U$ for every $t < s$, and $s \in [a, b] \subseteq I$, so $s$ lies in $U$ or in $V$. If $s \in V$, then since $V$ is open, $(s - \varepsilon, s] \subseteq V$ for some $\varepsilon > 0$; but points of $(s - \varepsilon, s)$ also lie in $U$, which contradicts $I \cap U \cap V = \emptyset$. If $s \in U$, then $s \neq b$ (since $b \in V$), so $s < b$, and since $U$ is open, $[s, s + \varepsilon] \subseteq U$ for some small $\varepsilon > 0$ with $s + \varepsilon \leq b$. Then $[a, s + \varepsilon] \subseteq U$, which contradicts the definition of $s$.

    (2) If $f(X)$ were covered by open sets $U$ and $V$ as in the definition, then the preimages $f^{-1}(U)$ and $f^{-1}(V)$ would disconnect $X$ (they are open by continuity, in the sense that they are of the form $X \cap U'$ for open sets $U'$).

    (3) Suppose that open sets $U$, $V$ disconnect a path-connected set $X$, with $x \in U$ and $y \in V$. Then the preimages of $U$ and $V$ under a path from $x$ to $y$ would disconnect the interval $[0, 1]$, which contradicts (1).
\end{proof}

Together, (1) and (2) give the intermediate value theorem: a continuous function $f: [a, b] \to \mathbb{R}$ takes every value between $f(a)$ and $f(b)$, since its image is a connected subset of $\mathbb{R}$, i.e., an interval.

\begin{example}\label{A-ex:connected}
    \begin{enumerate}
        \item $\mathbb{R}^n$, balls, and the sphere $S^{n-1}$ for $n \geq 2$ are path-connected, hence connected.
        \item The hyperbola $\{xy = 1\}$ is not connected: it lies in the union of the open half-planes $\{x > 0\}$ and $\{x < 0\}$, and meets both.
        \item $GL(n, \mathbb{R})$ is not connected: it is the union of the disjoint open sets $\{\det > 0\}$ and $\{\det < 0\}$. In particular, an invertible matrix with $\det < 0$ (a reflection) cannot be continuously deformed into the identity through invertible matrices. Similarly, $O(n)$ consists of the two pieces $\det Q = 1$ (the rotations, Lecture~2) and $\det Q = -1$; one can show that the rotations form a connected set.
        \item A circle and an interval are not homeomorphic. If we remove a point from a circle, the rest is still connected (it is an arc); if we remove an interior point from an interval, the rest is disconnected. A homeomorphism would map a point and its complement to a point and its complement, and it preserves connectedness by Proposition~\ref{A-prop:connected}(2).
    \end{enumerate}
\end{example}

\section{The Inverse Function Theorem}\label{A-sec:inverse}

From now on, we combine topology with derivatives. A map is smooth if all its partial derivatives of all orders exist and are continuous, and a diffeomorphism is a smooth bijection with a smooth inverse. We first fill a gap from Lecture~3, using a vector-valued version of the mean value theorem.

\begin{lemma}[Mean value inequality]\label{A-lem:mean-value}
    Let $\mathbf{g}: U \to \mathbb{R}^m$ be continuously differentiable on an open set $U \subseteq \mathbb{R}^n$ that contains the segment between $\mathbf{x}$ and $\mathbf{y}$. Then
    \[\norm{\mathbf{g}(\mathbf{y}) - \mathbf{g}(\mathbf{x})} \leq \sup_{\mathbf{z} \in [\mathbf{x}, \mathbf{y}]} \norm{D\mathbf{g}(\mathbf{z})}_2 \, \norm{\mathbf{y} - \mathbf{x}},\]
    where $\norm{\cdot}_2$ is the spectral norm of Lecture~2.
\end{lemma}

\begin{proof}
    Let $\mathbf{u} = \mathbf{g}(\mathbf{y}) - \mathbf{g}(\mathbf{x})$, and assume $\mathbf{u} \neq \mathbf{0}$. The function $\psi(t) = \mathbf{u}^\top\mathbf{g}(\mathbf{x} + t(\mathbf{y} - \mathbf{x}))$ has $\psi(1) - \psi(0) = \norm{\mathbf{u}}^2$. By the mean value theorem and the chain rule, $\psi(1) - \psi(0) = \psi'(\tau) = \mathbf{u}^\top D\mathbf{g}(\mathbf{z})(\mathbf{y} - \mathbf{x})$ with $\mathbf{z} = \mathbf{x} + \tau(\mathbf{y} - \mathbf{x})$. By the Cauchy--Schwarz inequality, $\norm{\mathbf{u}}^2 \leq \norm{\mathbf{u}}\,\norm{D\mathbf{g}(\mathbf{z})}_2\,\norm{\mathbf{y} - \mathbf{x}}$; divide by $\norm{\mathbf{u}}$.
\end{proof}

\begin{proposition}[Theorem~\ref{L3-thm:C1}]\label{A-prop:C1}
    If all partial derivatives of $f: U \to \mathbb{R}$ exist in a ball around $\mathbf{x}$ and are continuous at $\mathbf{x}$, then $f$ is differentiable at $\mathbf{x}$.
\end{proposition}

\begin{proof}
    For small $\mathbf{h}$, let $\mathbf{z}_0 = \mathbf{x}$ and $\mathbf{z}_j = \mathbf{x} + h_1\mathbf{e}_1 + \cdots + h_j\mathbf{e}_j$, so that $\mathbf{z}_n = \mathbf{x} + \mathbf{h}$. The points $\mathbf{z}_{j-1}$ and $\mathbf{z}_j$ differ only in the $j$-th coordinate, so by the mean value theorem in one variable, $f(\mathbf{z}_j) - f(\mathbf{z}_{j-1}) = \frac{\partial f}{\partial x_j}(\boldsymbol{\xi}_j)h_j$ for a point $\boldsymbol{\xi}_j$ between them. Adding these equations,
    \[\big|f(\mathbf{x} + \mathbf{h}) - f(\mathbf{x}) - \nabla f(\mathbf{x})^\top\mathbf{h}\big| = \Big|\sum_{j=1}^n \Big(\frac{\partial f}{\partial x_j}(\boldsymbol{\xi}_j) - \frac{\partial f}{\partial x_j}(\mathbf{x})\Big)h_j\Big| \leq \sum_{j=1}^n \Big|\frac{\partial f}{\partial x_j}(\boldsymbol{\xi}_j) - \frac{\partial f}{\partial x_j}(\mathbf{x})\Big| \, \norm{\mathbf{h}}.\]
    As $\mathbf{h} \to \mathbf{0}$, all $\boldsymbol{\xi}_j \to \mathbf{x}$, so the sum tends to $0$ by the continuity of the partial derivatives at $\mathbf{x}$.
\end{proof}

The key tool of this section is a simple but powerful theorem about solving equations by iteration.

\begin{theorem}[Banach fixed-point theorem]\label{A-thm:banach}
    Let $C \subseteq \mathbb{R}^n$ be closed and nonempty, and let $T: C \to C$ be a contraction, i.e., $\norm{T(\mathbf{x}) - T(\mathbf{y})} \leq q\norm{\mathbf{x} - \mathbf{y}}$ for all $\mathbf{x}, \mathbf{y} \in C$, with a constant $q < 1$. Then $T$ has exactly one fixed point $\mathbf{x}^* = T(\mathbf{x}^*)$. For every $\mathbf{x}_0 \in C$, the iterates $\mathbf{x}_{k+1} = T(\mathbf{x}_k)$ converge to $\mathbf{x}^*$, and
    \[\norm{\mathbf{x}_k - \mathbf{x}^*} \leq \frac{q^k}{1 - q}\norm{\mathbf{x}_1 - \mathbf{x}_0}.\]
\end{theorem}

\begin{proof}
    By induction, $\norm{\mathbf{x}_{j+1} - \mathbf{x}_j} \leq q^j\norm{\mathbf{x}_1 - \mathbf{x}_0}$. For $l > k$, the triangle inequality and the geometric series give
    \[\norm{\mathbf{x}_l - \mathbf{x}_k} \leq \sum_{j=k}^{l-1} q^j\norm{\mathbf{x}_1 - \mathbf{x}_0} \leq \frac{q^k}{1 - q}\norm{\mathbf{x}_1 - \mathbf{x}_0}.\]
    So $(\mathbf{x}_k)$ is a Cauchy sequence, which converges to some $\mathbf{x}^*$ by Corollary~\ref{A-cor:complete}, and $\mathbf{x}^* \in C$ since $C$ is closed. Since $T$ is continuous, $T(\mathbf{x}^*) = \lim T(\mathbf{x}_k) = \lim \mathbf{x}_{k+1} = \mathbf{x}^*$. Letting $l \to \infty$ in the inequality gives the error bound. If $\mathbf{y}^*$ is another fixed point, then $\norm{\mathbf{x}^* - \mathbf{y}^*} = \norm{T(\mathbf{x}^*) - T(\mathbf{y}^*)} \leq q\norm{\mathbf{x}^* - \mathbf{y}^*}$, so $\mathbf{x}^* = \mathbf{y}^*$.
\end{proof}

\begin{example}\label{A-ex:dottie}
    The function $T(x) = \cos x$ maps $[0, 1]$ into $[\cos 1, 1] \subseteq [0, 1]$, and $|T'(x)| = |\sin x| \leq \sin 1 \approx 0.841$ there, so it is a contraction with $q = \sin 1$ (by the mean value theorem). So the equation $\cos x = x$ has exactly one solution in $[0, 1]$, $x^* = 0.739085\ldots$, and pressing the cosine key of a calculator repeatedly finds it. Starting at $x_0 = 1$, the error bound guarantees an accuracy of $10^{-6}$ after $87$ steps; in fact, $32$ steps suffice. Gradient descent on a quadratic function with $\rho < 1$ (Theorem~\ref{L3-thm:gd-quadratic}) is another example of a contraction.
\end{example}

\begin{corollary}[Invertible residual layers]\label{A-cor:residual}
    Let $\mathbf{g}: \mathbb{R}^n \to \mathbb{R}^n$ satisfy $\norm{\mathbf{g}(\mathbf{x}) - \mathbf{g}(\mathbf{y})} \leq q\norm{\mathbf{x} - \mathbf{y}}$ with $q < 1$. Then the residual layer $\mathbf{f}(\mathbf{x}) = \mathbf{x} + \mathbf{g}(\mathbf{x})$ is a homeomorphism of $\mathbb{R}^n$, and $\mathbf{f}^{-1}(\mathbf{y})$ is the limit of the iteration $\mathbf{x}_{k+1} = \mathbf{y} - \mathbf{g}(\mathbf{x}_k)$.
\end{corollary}

\begin{proof}
    For fixed $\mathbf{y}$, the map $T(\mathbf{x}) = \mathbf{y} - \mathbf{g}(\mathbf{x})$ is a contraction of $\mathbb{R}^n$, and $\mathbf{f}(\mathbf{x}) = \mathbf{y}$ if and only if $\mathbf{x} = T(\mathbf{x})$. By Theorem~\ref{A-thm:banach}, there is exactly one solution, so $\mathbf{f}$ is bijective. It is continuous, and so is its inverse: $\norm{\mathbf{x} - \mathbf{x}'} \leq \norm{\mathbf{f}(\mathbf{x}) - \mathbf{f}(\mathbf{x}')} + \norm{\mathbf{g}(\mathbf{x}) - \mathbf{g}(\mathbf{x}')} \leq \norm{\mathbf{f}(\mathbf{x}) - \mathbf{f}(\mathbf{x}')} + q\norm{\mathbf{x} - \mathbf{x}'}$, so $\norm{\mathbf{f}^{-1}(\mathbf{y}) - \mathbf{f}^{-1}(\mathbf{y}')} \leq \frac{1}{1 - q}\norm{\mathbf{y} - \mathbf{y}'}$.
\end{proof}

Invertible neural networks used in normalizing flows (Remark~\ref{L3-rem:change-of-variables}) are built from such layers: constraining the Lipschitz constant of each residual branch (e.g., by normalizing the spectral norms of its weight matrices) makes the whole network invertible. A diffeomorphism must have invertible Jacobians everywhere, by the chain rule applied to $\mathbf{f}^{-1} \circ \mathbf{f} = \mathrm{id}$. The inverse function theorem is a local converse.

\begin{theorem}[Inverse function theorem]\label{A-thm:inverse-function}
    Let $\mathbf{f}: U \to \mathbb{R}^n$ be continuously differentiable on an open set $U \subseteq \mathbb{R}^n$, and let $D\mathbf{f}(\mathbf{a})$ be invertible. Then there are open sets $V \ni \mathbf{a}$ and $W \ni \mathbf{f}(\mathbf{a})$ such that $\mathbf{f}$ maps $V$ bijectively onto $W$, and the inverse $\mathbf{f}^{-1}: W \to V$ is continuously differentiable, with
    \[D(\mathbf{f}^{-1})(\mathbf{f}(\mathbf{x})) = D\mathbf{f}(\mathbf{x})^{-1} \qquad \text{for } \mathbf{x} \in V.\]
    If $\mathbf{f}$ is smooth, so is $\mathbf{f}^{-1}$.
\end{theorem}

\begin{proof}[Sketch of proof]
    Replacing $\mathbf{f}(\mathbf{x})$ by $D\mathbf{f}(\mathbf{a})^{-1}\big(\mathbf{f}(\mathbf{a} + \mathbf{x}) - \mathbf{f}(\mathbf{a})\big)$, we may assume $\mathbf{a} = \mathbf{0}$, $\mathbf{f}(\mathbf{0}) = \mathbf{0}$ and $D\mathbf{f}(\mathbf{0}) = I$. Then $\mathbf{g}(\mathbf{x}) = \mathbf{x} - \mathbf{f}(\mathbf{x})$ has $D\mathbf{g}(\mathbf{0}) = 0$, so by continuity there is an $r > 0$ such that the closed ball $\overline{B}_r = \{\norm{\mathbf{x}} \leq r\}$ lies in $U$ and $\norm{D\mathbf{g}(\mathbf{x})}_2 \leq \frac{1}{2}$ on it. By Lemma~\ref{A-lem:mean-value}, $\mathbf{g}$ is a contraction on $\overline{B}_r$ with $q = \frac{1}{2}$, and $\norm{\mathbf{g}(\mathbf{x})} \leq \frac{1}{2}\norm{\mathbf{x}}$. For $\norm{\mathbf{y}} < \frac{r}{2}$, the map $T(\mathbf{x}) = \mathbf{y} + \mathbf{g}(\mathbf{x})$ maps $\overline{B}_r$ into itself, since $\norm{T(\mathbf{x})} \leq \frac{r}{2} + \frac{r}{2}$, and it is a contraction. Its unique fixed point is the unique $\mathbf{x} \in \overline{B}_r$ with $\mathbf{f}(\mathbf{x}) = \mathbf{y}$, and it satisfies $\norm{\mathbf{x}} \leq \norm{\mathbf{y}} + \frac{1}{2}\norm{\mathbf{x}}$, i.e., $\norm{\mathbf{x}} \leq 2\norm{\mathbf{y}} < r$. So with $W = B(\mathbf{0}, \frac{r}{2})$ and $V = B(\mathbf{0}, r) \cap \mathbf{f}^{-1}(W)$, $\mathbf{f}: V \to W$ is bijective, and as in Corollary~\ref{A-cor:residual}, the inverse is continuous. One can show that it is also differentiable; then the chain rule applied to $\mathbf{f}^{-1}(\mathbf{f}(\mathbf{x})) = \mathbf{x}$ gives $D(\mathbf{f}^{-1})(\mathbf{f}(\mathbf{x}))\,D\mathbf{f}(\mathbf{x}) = I$.
\end{proof}

\begin{example}\label{A-ex:local-inverse}
    \begin{enumerate}
        \item Polar coordinates $\mathbf{f}(r, \varphi) = (r\cos\varphi, r\sin\varphi)$ have $\det D\mathbf{f} = r$ (Example~\ref{L3-ex:polar}). So near every point with $r > 0$, they have a smooth local inverse. But $\mathbf{f}$ is not injective, since $\mathbf{f}(r, \varphi + 2\pi) = \mathbf{f}(r, \varphi)$: the theorem is only local.
        \item The map $F(x, y) = (x^2 - y^2, 2xy)$, i.e., $z \mapsto z^2$ for $z = x + iy$, has $\det DF = 4(x^2 + y^2)$. Near every point $z \neq 0$, it has a local inverse, a ``branch of the square root''. Globally, it is two-to-one, since $F(-z) = F(z)$, and at $z = 0$, it is not even locally injective.
    \end{enumerate}
\end{example}

\section{The Implicit Function Theorem}\label{A-sec:implicit}

When can an equation $F(\mathbf{x}, \mathbf{y}) = \mathbf{0}$ be solved for some of the variables? We write points of $\mathbb{R}^{k + m}$ as $(\mathbf{x}, \mathbf{y})$ with $\mathbf{x} \in \mathbb{R}^k$ and $\mathbf{y} \in \mathbb{R}^m$, and split the Jacobian of $F: \mathbb{R}^{k+m} \to \mathbb{R}^m$ into the $m \times k$ matrix $D_{\mathbf{x}}F$ and the square $m \times m$ matrix $D_{\mathbf{y}}F$ of partial derivatives with respect to $\mathbf{x}$ and $\mathbf{y}$.

\begin{theorem}[Implicit function theorem]\label{A-thm:implicit-function}
    Let $F: U \to \mathbb{R}^m$ be continuously differentiable on an open set $U \subseteq \mathbb{R}^{k+m}$, let $F(\mathbf{a}, \mathbf{b}) = \mathbf{0}$, and let $D_{\mathbf{y}}F(\mathbf{a}, \mathbf{b})$ be invertible. Then there are open sets $V \ni \mathbf{a}$ and $W \ni (\mathbf{a}, \mathbf{b})$ and a continuously differentiable function $\mathbf{g}: V \to \mathbb{R}^m$ with $\mathbf{g}(\mathbf{a}) = \mathbf{b}$, such that for $(\mathbf{x}, \mathbf{y}) \in W$,
    \[F(\mathbf{x}, \mathbf{y}) = \mathbf{0} \quad \Longleftrightarrow \quad \mathbf{x} \in V \text{ and } \mathbf{y} = \mathbf{g}(\mathbf{x}).\]
    Moreover, $D\mathbf{g}(\mathbf{x}) = -D_{\mathbf{y}}F(\mathbf{x}, \mathbf{g}(\mathbf{x}))^{-1}\,D_{\mathbf{x}}F(\mathbf{x}, \mathbf{g}(\mathbf{x}))$.
\end{theorem}

\begin{proof}[Sketch of proof]
    Let $\Phi(\mathbf{x}, \mathbf{y}) = (\mathbf{x}, F(\mathbf{x}, \mathbf{y}))$. Its Jacobian $\begin{pmatrix} I & 0 \\ D_{\mathbf{x}}F & D_{\mathbf{y}}F \end{pmatrix}$ is invertible at $(\mathbf{a}, \mathbf{b})$: the equations $\mathbf{h} = \mathbf{u}$, $D_{\mathbf{x}}F\,\mathbf{h} + D_{\mathbf{y}}F\,\mathbf{k} = \mathbf{v}$ have the unique solution $\mathbf{h} = \mathbf{u}$, $\mathbf{k} = D_{\mathbf{y}}F^{-1}(\mathbf{v} - D_{\mathbf{x}}F\,\mathbf{u})$. By Theorem~\ref{A-thm:inverse-function}, $\Phi$ has a local inverse near $\Phi(\mathbf{a}, \mathbf{b}) = (\mathbf{a}, \mathbf{0})$. Since $\Phi$ does not change $\mathbf{x}$, neither does its inverse, which therefore has the form $(\mathbf{x}, \mathbf{z}) \mapsto (\mathbf{x}, G(\mathbf{x}, \mathbf{z}))$. Near $(\mathbf{a}, \mathbf{b})$, $F(\mathbf{x}, \mathbf{y}) = \mathbf{0}$ means $\Phi(\mathbf{x}, \mathbf{y}) = (\mathbf{x}, \mathbf{0})$, i.e., $\mathbf{y} = G(\mathbf{x}, \mathbf{0}) =: \mathbf{g}(\mathbf{x})$. Differentiating $F(\mathbf{x}, \mathbf{g}(\mathbf{x})) = \mathbf{0}$ with the chain rule gives $D_{\mathbf{x}}F + D_{\mathbf{y}}F\,D\mathbf{g} = 0$.
\end{proof}

\begin{example}\label{A-ex:implicit}
    \begin{enumerate}
        \item The unit circle $F(x, y) = x^2 + y^2 - 1 = 0$ has $\frac{\partial F}{\partial y} = 2y$. Near every point with $y \neq 0$, the circle is the graph of a function of $x$ ($y = \pm\sqrt{1 - x^2}$), with $\frac{dy}{dx} = -\frac{2x}{2y}$. Near $(\pm 1, 0)$, it is not, but it is the graph of a function of $y$, since $\frac{\partial F}{\partial x} = 2x \neq 0$ there.
        \item For the folium of Descartes $x^3 + y^3 = 3xy$ (Example~\ref{L3-ex:critical-points}), $\frac{\partial F}{\partial y} = 3y^2 - 3x$ vanishes on the curve only at $(0, 0)$, where the curve crosses itself, and at $(2^{2/3}, 2^{1/3})$, where its tangent is vertical. Near every other point of the curve, $y$ is a smooth function of $x$.
    \end{enumerate}
\end{example}

\begin{remark}[Implicit differentiation in machine learning]\label{A-rem:implicit-ml}
    Sometimes a quantity is defined as the solution of an equation, $G(\mathbf{z}, \boldsymbol{\theta}) = \mathbf{0}$, and we need its derivative with respect to $\boldsymbol{\theta}$. Examples are the minimizer $\mathbf{z}^*(\boldsymbol{\theta})$ of an inner optimization problem with hyperparameters $\boldsymbol{\theta}$ (here $G = \nabla_{\mathbf{z}}$ of the inner objective), and the fixed point $\mathbf{z} = \mathbf{h}(\mathbf{z}, \boldsymbol{\theta})$ of an ``implicit layer'' in deep equilibrium models. By Theorem~\ref{A-thm:implicit-function}, $D\mathbf{z}^*(\boldsymbol{\theta}) = -D_{\mathbf{z}}G^{-1}D_{\boldsymbol{\theta}}G$, so we do not need to backpropagate through all the iterations of the solver that computed $\mathbf{z}^*$. For example, the ridge regression solution $\mathbf{w}^*(\lambda)$ solves $G(\mathbf{w}, \lambda) = X^\top(X\mathbf{w} - \mathbf{y}) + \lambda\mathbf{w} = \mathbf{0}$, so
    \[\frac{d\mathbf{w}^*}{d\lambda} = -(X^\top X + \lambda I)^{-1}\mathbf{w}^*(\lambda),\]
    which can be used to tune $\lambda$ by gradient descent on a validation loss.
\end{remark}

\section{Submanifolds and Tangent Spaces}\label{A-sec:manifolds}

A $k$-dimensional manifold is a set that looks like $\mathbb{R}^k$ near each of its points: a curve looks like a line, and a surface like a plane. We consider manifolds that lie in $\mathbb{R}^n$ and are described by equations.

\begin{definition}[Submanifold]\label{A-def:submanifold}
    A set $M \subseteq \mathbb{R}^n$ is a $k$-dimensional submanifold of $\mathbb{R}^n$ if for every $\mathbf{p} \in M$, there are an open set $U \ni \mathbf{p}$ and a smooth map $F: U \to \mathbb{R}^{n-k}$ such that
    \[M \cap U = \{\mathbf{x} \in U : F(\mathbf{x}) = \mathbf{0}\} \qquad \text{and} \qquad \rank DF(\mathbf{p}) = n - k.\]
    (For $k = n$, there is no equation: $M$ is an open subset of $\mathbb{R}^n$.)
\end{definition}

So locally, $M$ is cut out by $n - k$ independent equations. In particular, if $F: U \to \mathbb{R}^m$ is smooth and $\mathbf{c} \in \mathbb{R}^m$ is a regular value of $F$, i.e., $\rank DF(\mathbf{x}) = m$ for every $\mathbf{x}$ with $F(\mathbf{x}) = \mathbf{c}$, then the level set $F^{-1}(\mathbf{c})$ is an $(n - m)$-dimensional submanifold (use $F - \mathbf{c}$ in the definition). The implicit function theorem shows that submanifolds really look like $\mathbb{R}^k$.

\begin{proposition}[Submanifolds are local graphs]\label{A-prop:local-graph}
    Let $M$ be a $k$-dimensional submanifold and $\mathbf{p} \in M$. After renumbering the coordinates of $\mathbb{R}^n = \mathbb{R}^k \times \mathbb{R}^{n-k}$, there are open sets $V \subseteq \mathbb{R}^k$ and $W \ni \mathbf{p}$ and a smooth map $\mathbf{g}: V \to \mathbb{R}^{n-k}$ such that $M \cap W = \{(\mathbf{x}, \mathbf{g}(\mathbf{x})) : \mathbf{x} \in V\}$. In particular, $\mathbf{x} \mapsto (\mathbf{x}, \mathbf{g}(\mathbf{x}))$ is a homeomorphism from the open set $V \subseteq \mathbb{R}^k$ onto $M \cap W$.
\end{proposition}

\begin{proof}
    The $(n - k) \times n$ matrix $DF(\mathbf{p})$ has rank $n - k$, so it has $n - k$ linearly independent columns (Lecture~1). Renumber the coordinates so that these are the last $n - k$ columns, and write $\mathbf{x} \in \mathbb{R}^n$ as $(\mathbf{x}', \mathbf{y})$. Then $D_{\mathbf{y}}F(\mathbf{p})$ is invertible, and Theorem~\ref{A-thm:implicit-function} gives $\mathbf{g}$. The inverse of $\mathbf{x}' \mapsto (\mathbf{x}', \mathbf{g}(\mathbf{x}'))$ is the projection to the first $k$ coordinates, which is continuous.
\end{proof}

\begin{example}\label{A-ex:submanifolds}
    \begin{enumerate}
        \item The sphere $S^{n-1} = \{\mathbf{x} : \norm{\mathbf{x}}^2 = 1\}$ is an $(n - 1)$-dimensional submanifold, since $F(\mathbf{x}) = \norm{\mathbf{x}}^2 - 1$ has $DF(\mathbf{x}) = 2\mathbf{x}^\top \neq \mathbf{0}$ on the sphere.
        \item The graph $\{(\mathbf{x}, \mathbf{g}(\mathbf{x}))\}$ of a smooth $\mathbf{g}: \mathbb{R}^k \to \mathbb{R}^{n-k}$ is a $k$-dimensional submanifold: $F(\mathbf{x}, \mathbf{y}) = \mathbf{y} - \mathbf{g}(\mathbf{x})$ has $DF = (-D\mathbf{g}, \, I)$, of rank $n - k$.
        \item $GL(n, \mathbb{R})$ is open in $\mathbb{R}^{n \times n}$, hence an $n^2$-dimensional submanifold.
        \item The special linear group $SL(n, \mathbb{R}) = \{X : \det X = 1\}$ is a submanifold of dimension $n^2 - 1$. By Jacobi's formula (Proposition~\ref{L3-prop:det-gradient}), $D\det(X)H = \det(X)\tr(X^{-1}H)$, which for $\det X = 1$ and $H = X$ equals $n \neq 0$. So $D\det(X) \neq 0$, and $1$ is a regular value of $\det$.
        \item The orthogonal group $O(n) = \{Q : Q^\top Q = I\}$ is a submanifold of dimension $\frac{n(n-1)}{2}$. Let $F(Q) = Q^\top Q$, a map from $\mathbb{R}^{n \times n}$ to the space $\mathrm{Sym}(n)$ of symmetric matrices, which has dimension $\frac{n(n+1)}{2}$. Its derivative is $DF(Q)H = H^\top Q + Q^\top H$. For $Q \in O(n)$ and any symmetric $S$, the matrix $H = \frac{1}{2}QS$ satisfies $DF(Q)H = \frac{1}{2}SQ^\top Q + \frac{1}{2}Q^\top QS = S$, so $DF(Q)$ is surjective, i.e., has full rank $\frac{n(n+1)}{2}$, and $O(n)$ has dimension $n^2 - \frac{n(n+1)}{2} = \frac{n(n-1)}{2}$. For $n = 2$, this is $1$ (one rotation angle), and for $n = 3$, it is $3$ (three angles, as for the rotations $R_1$, $R_2$, $R_3$ of Lecture~2).
    \end{enumerate}
\end{example}

Not every level set is a submanifold. The level set $\{x^2 - y^2 = 0\}$ of Figure~\ref{L3-fig:level-sets} in Lecture~3 consists of two lines that cross at the origin, where the gradient $(2x, -2y)$ vanishes. It is not a $1$-dimensional submanifold near the origin: removing the origin from a small piece of it leaves four connected pieces, while removing a point from a small piece of a curve that looks like a line leaves two. Similarly, the double cone $\{x^2 + y^2 = z^2\}$ is not a submanifold at its tip. On the other hand, the rank condition is sufficient, but not necessary: $\{x^2 = 0\}$ is the $y$-axis, a perfectly good submanifold, although $D(x^2) = (2x, 0)$ vanishes on it.

\begin{definition}[Tangent space]\label{A-def:tangent}
    Let $M$ be a submanifold and $\mathbf{p} \in M$. The tangent space $T_{\mathbf{p}}M$ is the set of velocities $\boldsymbol{\gamma}'(0)$ of all smooth curves $\boldsymbol{\gamma}: (-\varepsilon, \varepsilon) \to \mathbb{R}^n$ with $\boldsymbol{\gamma}(t) \in M$ for all $t$ and $\boldsymbol{\gamma}(0) = \mathbf{p}$. Its orthogonal complement $N_{\mathbf{p}}M = (T_{\mathbf{p}}M)^\perp$ is the normal space.
\end{definition}

\begin{theorem}\label{A-thm:tangent-space}
    Let $M$ be a $k$-dimensional submanifold, and let $M \cap U = F^{-1}(\mathbf{0})$ near $\mathbf{p}$ as in Definition~\ref{A-def:submanifold}. Then
    \[T_{\mathbf{p}}M = \Null DF(\mathbf{p}) \qquad \text{and} \qquad N_{\mathbf{p}}M = \row DF(\mathbf{p}) = \Span\{\nabla F_1(\mathbf{p}), \ldots, \nabla F_{n-k}(\mathbf{p})\}.\]
    In particular, $T_{\mathbf{p}}M$ is a $k$-dimensional subspace.
\end{theorem}

\begin{proof}
    If $\boldsymbol{\gamma}$ is a curve in $M$ through $\mathbf{p}$, then $F(\boldsymbol{\gamma}(t)) = \mathbf{0}$ for small $t$, and the chain rule gives $DF(\mathbf{p})\boldsymbol{\gamma}'(0) = \mathbf{0}$. So $T_{\mathbf{p}}M \subseteq \Null DF(\mathbf{p})$, which has dimension $n - (n - k) = k$ by the rank--nullity theorem (Lecture~1). Conversely, write $M$ near $\mathbf{p} = (\mathbf{a}, \mathbf{g}(\mathbf{a}))$ as a graph as in Proposition~\ref{A-prop:local-graph}. For every $\mathbf{v} \in \mathbb{R}^k$, the curve $t \mapsto (\mathbf{a} + t\mathbf{v}, \mathbf{g}(\mathbf{a} + t\mathbf{v}))$ lies in $M$ and has velocity $(\mathbf{v}, D\mathbf{g}(\mathbf{a})\mathbf{v})$ at $t = 0$. These velocities form a $k$-dimensional subspace (the image of an injective linear map), so $T_{\mathbf{p}}M$ contains a $k$-dimensional subspace of the $k$-dimensional space $\Null DF(\mathbf{p})$, and the two are equal. Finally, $\Null(A)^\perp = \row(A)$ for every matrix $A$ (Lecture~2), and the rows of $DF(\mathbf{p})$ are the transposed gradients $\nabla F_i(\mathbf{p})^\top$.
\end{proof}

\begin{figure}[ht]
    \centering
    \begin{tikzpicture}[scale=1.5, >=Stealth]
        \draw[->, gray] (-1.5,0) -- (1.7,0) node[right] {$x$};
        \draw[->, gray] (0,-1.3) -- (0,1.5) node[above] {$y$};
        \draw[blue!70!black, thick] (0,0) circle (1);
        \coordinate (P) at (30:1);
        \draw[orange!85!black, thick] ($(P) + (120:1.1)$) -- ($(P) + (-60:1.1)$);
        \draw[->, very thick, red!75!black] (P) -- ($(P) + (30:0.7)$) node[right] {$\nabla F(\mathbf{p}) = 2\mathbf{p}$};
        \draw[->, very thick, orange!85!black] (P) -- ($(P) + (120:0.6)$) node[above, fill=white, inner sep=1pt] {$\mathbf{v} \in T_{\mathbf{p}}M$};
        \fill (P) circle (1.2pt) node[below left] {$\mathbf{p}$};
        \node[blue!70!black] at (-0.95,-0.95) {$M$};
    \end{tikzpicture}
    \caption{The circle $M = \{F = 0\}$ with $F(x, y) = x^2 + y^2 - 1$. The tangent space $T_{\mathbf{p}}M = \Null DF(\mathbf{p}) = \mathbf{p}^\perp$ is the tangent line (shifted to $\mathbf{p}$), and the normal space is spanned by the gradient $\nabla F(\mathbf{p})$.}
    \label{A-fig:tangent}
\end{figure}

\begin{example}\label{A-ex:tangent-spaces}
    \begin{enumerate}
        \item $T_{\mathbf{p}}S^{n-1} = \Null(2\mathbf{p}^\top) = \{\mathbf{v} : \mathbf{p}^\top\mathbf{v} = 0\} = \mathbf{p}^\perp$: the tangent vectors of a sphere are orthogonal to the radius (Figure~\ref{A-fig:tangent}).
        \item $T_I SL(n, \mathbb{R}) = \{H : D\det(I)H = 0\} = \{H : \tr H = 0\}$, the traceless matrices: to first order, $\det(I + tH) = 1 + t\tr H$ (Lemma~\ref{L3-lem:det-identity}), so a matrix curve through $I$ keeps determinant $1$ only if its velocity has trace $0$.
        \item $T_Q O(n) = \{H : H^\top Q + Q^\top H = 0\}$. With $K = Q^\top H$, the condition says $K^\top + K = 0$, so $T_Q O(n) = \{QK : K^\top = -K\}$, and $T_I O(n)$ is the space of skew-symmetric matrices, of dimension $\frac{n(n-1)}{2}$. For example, the rotations $R(t)$ by the angle $t$ in the plane have $R'(0) = \begin{pmatrix} 0 & -1 \\ 1 & 0 \end{pmatrix}$.
    \end{enumerate}
\end{example}

\section{Lagrange Multipliers and Optimization on Manifolds}\label{A-sec:lagrange}

We can now complete the proof of the Lagrange multiplier theorem (Theorem~\ref{L3-thm:lagrange}).

\begin{theorem}[Lagrange multipliers]\label{A-thm:lagrange-proof}
    Let $f$ and $g_1, \ldots, g_m$ be smooth on an open set $U \subseteq \mathbb{R}^n$, let $\mathbf{x}^*$ be a local minimum of $f$ on $M = \{\mathbf{x} \in U : g_1(\mathbf{x}) = \cdots = g_m(\mathbf{x}) = 0\}$, and let $\nabla g_1(\mathbf{x}^*), \ldots, \nabla g_m(\mathbf{x}^*)$ be linearly independent. Then $\nabla f(\mathbf{x}^*) \in \Span\{\nabla g_1(\mathbf{x}^*), \ldots, \nabla g_m(\mathbf{x}^*)\}$, i.e., there are $\lambda_i$ with $\nabla f(\mathbf{x}^*) + \sum_i \lambda_i\nabla g_i(\mathbf{x}^*) = \mathbf{0}$.
\end{theorem}

\begin{proof}
    The linear independence says that $DG(\mathbf{x}^*)$ has rank $m$ for $G = (g_1, \ldots, g_m)$, so near $\mathbf{x}^*$, $M$ is an $(n - m)$-dimensional submanifold (the rank stays $m$ near $\mathbf{x}^*$, since a nonzero $m \times m$ minor stays nonzero by continuity). For every $\mathbf{v} \in T_{\mathbf{x}^*}M$, choose a curve $\boldsymbol{\gamma}$ in $M$ with $\boldsymbol{\gamma}(0) = \mathbf{x}^*$ and $\boldsymbol{\gamma}'(0) = \mathbf{v}$. The function $t \mapsto f(\boldsymbol{\gamma}(t))$ has a local minimum at $t = 0$, so its derivative $\nabla f(\mathbf{x}^*)^\top\mathbf{v}$ is $0$. Hence $\nabla f(\mathbf{x}^*) \in (T_{\mathbf{x}^*}M)^\perp = N_{\mathbf{x}^*}M = \Span\{\nabla g_i(\mathbf{x}^*)\}$ by Theorem~\ref{A-thm:tangent-space}.
\end{proof}

The proof suggests how to do calculus on a manifold: only the component of $\nabla f$ in the tangent space matters.

\begin{definition}[Riemannian gradient]\label{A-def:riemannian-gradient}
    Let $M$ be a submanifold of $\mathbb{R}^n$, $\mathbf{p} \in M$, and $f$ smooth near $M$. The Riemannian gradient of $f$ on $M$ at $\mathbf{p}$ is the orthogonal projection (Lecture~2) of $\nabla f(\mathbf{p})$ onto $T_{\mathbf{p}}M$:
    \[\grad_M f(\mathbf{p}) = P_{\mathbf{p}}\nabla f(\mathbf{p}).\]
\end{definition}

Since $\nabla f(\mathbf{p}) - \grad_M f(\mathbf{p})$ is orthogonal to $T_{\mathbf{p}}M$, we have $\nabla f(\mathbf{p})^\top\mathbf{v} = \grad_M f(\mathbf{p})^\top\mathbf{v}$ for all tangent vectors $\mathbf{v}$: the Riemannian gradient gives the rates of change of $f$ along $M$. As in Theorem~\ref{L3-thm:steepest}, it points in the direction of steepest ascent among the tangent directions, and by Theorem~\ref{A-thm:lagrange-proof}, $\grad_M f(\mathbf{p}) = \mathbf{0}$ exactly at the points that satisfy the Lagrange condition. These points are the critical points of $f$ on $M$.

Riemannian gradient descent minimizes $f$ on $M$ by taking a step in the direction $-\grad_M f(\mathbf{p}_k)$, which leaves $M$ slightly, and mapping the result back to $M$ with a retraction $R$:
\[\mathbf{p}_{k+1} = R\big(\mathbf{p}_k - \eta\grad_M f(\mathbf{p}_k)\big).\]
On the sphere, $P_{\mathbf{p}} = I - \mathbf{p}\mathbf{p}^\top$, so $\grad_M f(\mathbf{p}) = \nabla f(\mathbf{p}) - (\mathbf{p}^\top\nabla f(\mathbf{p}))\mathbf{p}$, and the retraction can simply normalize: $R(\mathbf{x}) = \frac{\mathbf{x}}{\norm{\mathbf{x}}}$. On $O(n)$, one can retract with the $Q$ factor of the QR decomposition (Lecture~2). Such methods are used for PCA, for orthogonality constraints on weight matrices, and for low-rank matrix problems.

\begin{example}[The Rayleigh quotient on the sphere]\label{A-ex:rayleigh-sphere}
    Let $A$ be symmetric with eigenvalues $\lambda_1 > \lambda_2 > \cdots > \lambda_n$ and unit eigenvectors $\mathbf{q}_1, \ldots, \mathbf{q}_n$, and let $f(\mathbf{x}) = \mathbf{x}^\top A\mathbf{x}$ on $S^{n-1}$. Then $\grad_M f(\mathbf{x}) = 2A\mathbf{x} - 2(\mathbf{x}^\top A\mathbf{x})\mathbf{x}$, which is $\mathbf{0}$ if and only if $\mathbf{x}$ is an eigenvector: the critical points are the $2n$ points $\pm\mathbf{q}_i$. To see their types, move from $\mathbf{q}_i$ towards $\mathbf{q}_j$ along the great circle $\boldsymbol{\gamma}(t) = \cos t\,\mathbf{q}_i + \sin t\,\mathbf{q}_j$:
    \[f(\boldsymbol{\gamma}(t)) = \lambda_i\cos^2 t + \lambda_j\sin^2 t = \lambda_i + (\lambda_j - \lambda_i)\sin^2 t.\]
    So $f$ decreases in the $n - i$ tangent directions $\mathbf{q}_j$ with $\lambda_j < \lambda_i$, and increases in the $i - 1$ directions with $\lambda_j > \lambda_i$. Hence $\pm\mathbf{q}_1$ are the maxima, $\pm\mathbf{q}_n$ the minima, and all other eigenvectors are saddle points. In particular, Riemannian gradient ascent for the first principal component (Lecture~2) has exactly one pair of local maxima, the correct answer; the other critical points are saddle points, which gradient ascent avoids from almost every starting point.
\end{example}

\section{Critical Points and Morse Theory}\label{A-sec:morse}

\subsection{Regular and critical values}

\begin{definition}[Critical points and values of maps]\label{A-def:critical-maps}
    Let $F: U \to \mathbb{R}^m$ be smooth on an open set $U \subseteq \mathbb{R}^n$. A point $\mathbf{x}$ is a critical point of $F$ if $\rank DF(\mathbf{x}) < m$, and a value $\mathbf{c}$ is a critical value if $\mathbf{c} = F(\mathbf{x})$ for a critical point $\mathbf{x}$. All other values in $\mathbb{R}^m$ are regular values.
\end{definition}

For $m = 1$, the critical points are the points with $\nabla F(\mathbf{x}) = \mathbf{0}$, as in Lecture~3. By Section~\ref{A-sec:manifolds}, the level set of a regular value is a submanifold (or empty). How many critical values can there be? A set $Z \subseteq \mathbb{R}^m$ has measure zero if for every $\varepsilon > 0$ it can be covered by countably many boxes whose volumes add up to less than $\varepsilon$. For example, finite and countable sets, and lines in $\mathbb{R}^2$, have measure zero.

\begin{theorem}[Sard]\label{A-thm:sard}
    The set of critical values of a smooth map $F: U \to \mathbb{R}^m$ has measure zero.
\end{theorem}

We do not prove this theorem. It says that almost every value is regular, so almost every level set of a smooth function $f: U \to \mathbb{R}$ is an $(n - 1)$-dimensional submanifold or empty. Note that the theorem is about values, not points: a constant function has only critical points, but only one critical value.

\subsection{Nondegenerate critical points}

From now on, let $f: U \to \mathbb{R}$ be smooth. Lecture~3 classified critical points by the eigenvalues of the Hessian (Theorem~\ref{L3-thm:second-order}). The test was inconclusive when an eigenvalue is $0$.

\begin{definition}[Nondegenerate critical points and the index]\label{A-def:index}
    A critical point $\mathbf{p}$ of $f$ is nondegenerate if the Hessian $\nabla^2 f(\mathbf{p})$ is invertible, i.e., has no eigenvalue $0$. Its index is the number of negative eigenvalues of $\nabla^2 f(\mathbf{p})$, counted with multiplicity. A function whose critical points are all nondegenerate is a Morse function.
\end{definition}

So a nondegenerate critical point of index $0$ is a strict local minimum, one of index $n$ is a strict local maximum, and all others are saddle points. The index counts the independent directions in which $f$ decreases. For example, $x^3 + y^3 - 3xy$ (Example~\ref{L3-ex:critical-points}) is a Morse function, with a critical point of index $1$ at the origin and one of index $0$ at $(1, 1)$. The ``monkey saddle'' $x^3 - 3xy^2$ has a degenerate critical point at the origin, where its Hessian is $0$; near the origin, it goes up and down three times, unlike any nondegenerate saddle.

\begin{proposition}\label{A-prop:isolated}
    Nondegenerate critical points are isolated: each one has a neighborhood that contains no other critical point.
\end{proposition}

\begin{proof}
    Let $H = \nabla^2 f(\mathbf{p})$, and let $\mu > 0$ be the smallest absolute value of an eigenvalue of $H$. Writing $\mathbf{h}$ in an orthonormal eigenbasis (Lecture~2) shows $\norm{H\mathbf{h}} \geq \mu\norm{\mathbf{h}}$. The gradient is differentiable with derivative $H$ at $\mathbf{p}$, so $\nabla f(\mathbf{p} + \mathbf{h}) = H\mathbf{h} + \mathbf{r}(\mathbf{h})$ with $\norm{\mathbf{r}(\mathbf{h})} \leq \frac{\mu}{2}\norm{\mathbf{h}}$ for small $\mathbf{h}$. Then $\norm{\nabla f(\mathbf{p} + \mathbf{h})} \geq \frac{\mu}{2}\norm{\mathbf{h}} > 0$ for small $\mathbf{h} \neq \mathbf{0}$.
\end{proof}

In Lecture~3, we saw that near a critical point, $f$ is approximately a sum of parabolas, $f(\mathbf{p}) + \frac{1}{2}\sum_i \lambda_i y_i^2$, in the coordinates given by the eigenvectors of the Hessian. For nondegenerate critical points, this becomes exact after a nonlinear change of coordinates.

\begin{theorem}[Morse lemma]\label{A-thm:morse-lemma}
    Let $\mathbf{p}$ be a nondegenerate critical point of $f$ with index $k$. Then there is a diffeomorphism $\boldsymbol{\varphi}$ from a neighborhood of $\mathbf{p}$ onto a neighborhood of $\mathbf{0}$, with $\boldsymbol{\varphi}(\mathbf{p}) = \mathbf{0}$, such that in the coordinates $\mathbf{y} = \boldsymbol{\varphi}(\mathbf{x})$,
    \[f = f(\mathbf{p}) - y_1^2 - \cdots - y_k^2 + y_{k+1}^2 + \cdots + y_n^2.\]
\end{theorem}

We do not prove the Morse lemma. It says that near a nondegenerate critical point, the higher-order terms of the Taylor expansion do not change the picture at all: up to a smooth deformation, the function is a bowl, a saddle with $k$ downward directions, or an upside-down bowl. Degenerate critical points can be much more complicated, but they are rare.

\begin{theorem}[Morse functions are generic]\label{A-thm:morse-generic}
    Let $f: U \to \mathbb{R}$ be smooth. For almost every $\mathbf{a} \in \mathbb{R}^n$ (all except a set of measure zero), the function $f_{\mathbf{a}}(\mathbf{x}) = f(\mathbf{x}) - \mathbf{a}^\top\mathbf{x}$ is a Morse function.
\end{theorem}

\begin{proof}
    We have $\nabla f_{\mathbf{a}} = \nabla f - \mathbf{a}$ and $\nabla^2 f_{\mathbf{a}} = \nabla^2 f$. So $\mathbf{x}$ is a critical point of $f_{\mathbf{a}}$ if and only if $\nabla f(\mathbf{x}) = \mathbf{a}$, and it is degenerate if and only if $\nabla^2 f(\mathbf{x})$, the Jacobian of the map $\nabla f: U \to \mathbb{R}^n$, is singular. Hence $f_{\mathbf{a}}$ has a degenerate critical point if and only if $\mathbf{a}$ is a critical value of $\nabla f$, and by Sard's theorem, these $\mathbf{a}$ form a set of measure zero.
\end{proof}

For example, $f(x) = x^3$ has a degenerate critical point at $0$, but $x^3 - ax$ is a Morse function for every $a \neq 0$: it has two nondegenerate critical points if $a > 0$, and none if $a < 0$. An arbitrarily small, random linear perturbation removes all degenerate critical points.

\subsection{Critical points and the shape of sublevel sets}

Morse theory studies how the sublevel sets $M^c = \{\mathbf{x} : f(\mathbf{x}) \leq c\}$ change as $c$ increases. We state its two main results informally, for a Morse function $f$ whose sublevel sets are compact (for example, a Morse function on a compact manifold, or a coercive Morse function on $\mathbb{R}^n$).
\begin{itemize}
    \item If there are no critical values in $[a, b]$, then $M^a$ and $M^b$ have the same shape: $M^b$ can be continuously deformed into $M^a$, by flowing downhill along $-\nabla f$. Nothing interesting happens between critical values.
    \item When $c$ passes the value of a critical point of index $k$, the sublevel set changes by attaching a $k$-dimensional ``cell'' (a point for $k = 0$, a segment for $k = 1$, a disk for $k = 2$, and so on). A minimum ($k = 0$) creates a new connected component. A saddle point of index $1$ attaches a bridge, which either connects two components or closes a loop.
\end{itemize}

\begin{example}[A double well]\label{A-ex:double-well}
    The function $f(x, y) = (x^2 - 1)^2 + y^2$ has $\nabla f = (4x(x^2 - 1), 2y)$, so its critical points are $(\pm 1, 0)$, with $f = 0$ and $\nabla^2 f = \diag(8, 2)$ (index $0$), and $(0, 0)$, with $f = 1$ and $\nabla^2 f = \diag(-4, 2)$ (index $1$). For $c < 0$, the sublevel set is empty; at $c = 0$, two components appear at the minima; at $c = 1$, the saddle point attaches a bridge that joins them, and for $c > 1$, the sublevel set is connected (Figure~\ref{A-fig:double-well}). For gradient descent, this means that the two basins of the minima are separated by the saddle: an iteration that starts with $f(\mathbf{x}_0) < 1$ and decreases $f$ in every step cannot move from one basin to the other.
\end{example}

\begin{figure}[ht]
    \centering
    \foreach \c/\lab in {0.5/{c = 0.5}, 1/{c = 1}, 1.8/{c = 1.8}} {
        \begin{tikzpicture}[scale=1.15, >=Stealth]
            \pgfmathsetmacro{\s}{sqrt(\c)}
            \ifdim \c pt < 1.01pt
                \foreach \sgn in {1,-1}
                    \filldraw[fill=blue!15, draw=blue!70!black, thick, domain=-180:180, samples=120, variable=\t, smooth]
                        plot ({\sgn*sqrt(max(0, 1 + \s*cos(\t)))}, {\s*sin(\t)}) -- cycle;
            \else
                \pgfmathsetmacro{\tz}{acos(-1/\s)}
                \filldraw[fill=blue!15, draw=blue!70!black, thick, smooth]
                    plot[domain=-\tz:\tz, samples=90, variable=\t] ({sqrt(max(0, 1 + \s*cos(\t)))}, {\s*sin(\t)})
                    -- plot[domain=\tz:-\tz, samples=90, variable=\t] ({-sqrt(max(0, 1 + \s*cos(\t)))}, {\s*sin(\t)}) -- cycle;
            \fi
            \draw[->, gray] (-1.75,0) -- (1.85,0) node[right] {$x$};
            \draw[->, gray] (0,-1.45) -- (0,1.55) node[above] {$y$};
            \fill (1,0) circle (1.3pt);
            \fill (-1,0) circle (1.3pt);
            \fill[red!75!black] (0,0) circle (1.3pt);
            \node at (0,-1.75) {$\lab$};
        \end{tikzpicture}\hspace{0.6em}
    }
    \caption{Sublevel sets $\{f \leq c\}$ of the double well $f(x, y) = (x^2 - 1)^2 + y^2$, with the two minima (black, index $0$, value $0$) and the saddle point (red, index $1$, value $1$). Below the critical value $1$, the sublevel set has two components; at $c = 1$, they touch at the saddle point; above it, they have merged into one.}
    \label{A-fig:double-well}
\end{figure}

On a compact manifold, these changes must add up. The Euler characteristic $\chi(M)$ is a number that only depends on the topology of $M$: for a surface divided into $V$ vertices, $E$ edges and $F$ faces, it is $V - E + F$, which does not depend on the subdivision. For example, the sphere has $\chi = 2$ (a cube has $8 - 12 + 6 = 2$), and the torus (the surface of a doughnut) has $\chi = 0$.

\begin{theorem}[Morse relation]\label{A-thm:morse-relation}
    Let $f$ be a Morse function on a compact manifold $M$, and let $c_k$ be the number of its critical points of index $k$. Then
    \[\chi(M) = c_0 - c_1 + c_2 - \cdots + (-1)^n c_n.\]
\end{theorem}

We do not prove this theorem; it follows from the fact that attaching a $k$-cell changes the Euler characteristic by $(-1)^k$. It shows that the topology of $M$ forces critical points to exist, whatever function we choose.

\begin{example}\label{A-ex:morse-relation}
    \begin{enumerate}
        \item The height function on the sphere $S^2$ has a minimum and a maximum: $\chi(S^2) = 1 - 0 + 1 = 2$.
        \item The height function on an upright torus has a minimum, two saddle points and a maximum (Figure~\ref{A-fig:torus}): $\chi = 1 - 2 + 1 = 0$. In fact, every Morse function on the torus has at least four critical points, at least two of which are saddle points (this follows from the Morse inequalities, which we do not discuss).
        \item For the Rayleigh quotient on $S^{n-1}$ (Example~\ref{A-ex:rayleigh-sphere}), the two critical points $\pm\mathbf{q}_i$ have index $n - i$, so $\sum_k (-1)^k c_k = 2\sum_{i=1}^n (-1)^{n-i}$, which is $2$ for odd $n$ and $0$ for even $n$. This agrees with the Euler characteristic of the sphere $S^{n-1}$, which is $2$ if $n - 1$ is even and $0$ if $n - 1$ is odd. For $n = 3$ and $A = \diag(3, 2, 1)$: two maxima $\pm\mathbf{e}_1$, two saddle points $\pm\mathbf{e}_2$, two minima $\pm\mathbf{e}_3$, and $2 - 2 + 2 = 2$.
    \end{enumerate}
\end{example}

\begin{figure}[ht]
    \centering
    \begin{tikzpicture}[scale=0.95, >=Stealth]
        \begin{scope}[rotate=90]
            \draw[thick] (0,0) ellipse (2.4 and 1.35);
            \draw[thick] (-1.35,0.12) .. controls (-0.75,-0.62) and (0.75,-0.62) .. (1.35,0.12);
            \draw[thick] (-1.08,-0.18) .. controls (-0.55,0.35) and (0.55,0.35) .. (1.08,-0.18);
            \fill (-2.4,0) circle (1.6pt);
            \fill[red!75!black] (-1.15,-0.08) circle (1.6pt);
            \fill[red!75!black] (1.15,-0.08) circle (1.6pt);
            \fill (2.4,0) circle (1.6pt);
        \end{scope}
        \draw[dashed, gray] (0.15,-2.4) -- (3.0,-2.4) node[right, black] {minimum (index $0$)};
        \draw[dashed, gray] (0.15,-1.15) -- (3.0,-1.15) node[right, black] {saddle point (index $1$)};
        \draw[dashed, gray] (0.15,1.15) -- (3.0,1.15) node[right, black] {saddle point (index $1$)};
        \draw[dashed, gray] (0.15,2.4) -- (3.0,2.4) node[right, black] {maximum (index $2$)};
        \draw[->] (-2.4,-2.4) -- (-2.4,2.6) node[above] {height};
    \end{tikzpicture}
    \caption{The height function on an upright torus has four critical points: the bottom (a minimum), the bottom and the top of the hole (two saddle points), and the top (a maximum). As the height $c$ increases, the sublevel set is first a disk around the bottom (after the minimum), then a bent tube, i.e., a cylinder (after the first saddle point, where the disk closes up around the tube), then a torus with a hole at its top (after the second saddle point, where the two ends of the tube meet and close a loop around the hole of the torus), and finally the whole torus (after the maximum).}
    \label{A-fig:torus}
\end{figure}

\section{Topology in Machine Learning}\label{A-sec:ml}

We end with a few places where the ideas of this appendix appear in machine learning. Some of them are theorems, others are well-supported heuristics and active research topics.
\begin{itemize}
    \item \textbf{The manifold hypothesis.} High-dimensional data, such as images with millions of pixels, are believed to lie close to submanifolds of much smaller dimension: the images of a face, for example, vary with a small number of parameters such as pose, lighting and expression. PCA finds the best approximating linear subspace (Lecture~2); autoencoders and other nonlinear methods try to find curved, manifold-like structures.
    \item \textbf{Normalizing flows and topology.} A normalizing flow maps a simple distribution, such as a Gaussian on $\mathbb{R}^n$, to a complicated one by a diffeomorphism (Corollary~\ref{A-cor:residual} and Remark~\ref{L3-rem:change-of-variables}). Diffeomorphisms preserve topology. So, by Proposition~\ref{A-prop:connected}, a flow cannot map $\mathbb{R}^n$, which is connected, exactly onto a distribution supported on two separate pieces, such as two clusters: it can only approximate it by squeezing a little probability mass through thin bridges between them.
    \item \textbf{Symmetries make critical points degenerate.} If a loss is invariant under a continuous family of transformations of the parameters, its minima are not isolated. For example, the ReLU satisfies $\max(0, ct) = c\max(0, t)$ for $c > 0$, so multiplying the incoming weights of a hidden unit by $c$ and its outgoing weights by $\frac{1}{c}$ does not change the network. The set of minima then contains curves (in the simplest case, the hyperbola $w_1 w_2 = 1$ for the loss $(w_1 w_2 - 1)^2$), and by Proposition~\ref{A-prop:isolated}, the Hessian at a minimum must have the eigenvalue $0$, in the tangent directions of these curves. Loss functions of neural networks are therefore not Morse functions.
    \item \textbf{Saddle points in high dimensions.} At a critical point of a loss with millions of parameters, all eigenvalues of the Hessian must be nonnegative for a local minimum. A heuristic, supported by experiments and by random matrix theory, says that critical points with high loss usually have many negative eigenvalues, i.e., a large index: they are saddle points rather than bad local minima. Gradient descent with noise, as in SGD, tends to escape from saddle points, which is part of the explanation why training works.
    \item \textbf{Connected sublevel sets.} Experiments show that two minima found by training the same network twice can often be joined by a path along which the loss stays low (``mode connectivity''): the sublevel sets $\{L \leq c\}$ for small $c$ seem to be connected, up to the symmetries of the network. Morse theory gives a language for such statements: a sublevel set can only become connected by passing saddle points of index $1$.
\end{itemize}

\section*{Key Takeaways}
\addcontentsline{toc}{section}{Key Takeaways}
\begin{itemize}[leftmargin=*]
    \item Open sets, closed sets, convergence and continuity only depend on the notion of nearness. A map is continuous if and only if preimages of open sets are open; so sublevel sets $\{f \leq c\}$ and level sets of continuous functions are closed. All norms on $\mathbb{R}^n$ define the same open sets.
    \item A subset of $\mathbb{R}^n$ is compact if and only if it is closed and bounded (Heine--Borel). A continuous function on a nonempty compact set attains its minimum and maximum, and a continuous coercive function on $\mathbb{R}^n$ has a global minimum; regularization makes losses coercive.
    \item Continuous maps preserve connectedness; this gives the intermediate value theorem, and shows, e.g., that $GL(n, \mathbb{R})$ has two pieces and that a circle is not homeomorphic to an interval.
    \item A contraction of a closed set has a unique fixed point, found by iteration (Banach). This proves the inverse function theorem: a map with an invertible Jacobian is locally invertible, with $D(\mathbf{f}^{-1}) = (D\mathbf{f})^{-1}$. Residual layers $\mathbf{x} + \mathbf{g}(\mathbf{x})$ with a contractive $\mathbf{g}$ are invertible.
    \item The implicit function theorem solves $F(\mathbf{x}, \mathbf{y}) = \mathbf{0}$ locally for $\mathbf{y} = \mathbf{g}(\mathbf{x})$ when $D_{\mathbf{y}}F$ is invertible, with $D\mathbf{g} = -D_{\mathbf{y}}F^{-1}D_{\mathbf{x}}F$; this is implicit differentiation, also used for implicit layers and hyperparameter optimization.
    \item The level set of a regular value of $F: \mathbb{R}^n \to \mathbb{R}^m$ is an $(n - m)$-dimensional submanifold, which is locally a graph. Its tangent space is $\Null DF(\mathbf{p})$, and its normal space is spanned by the gradients $\nabla F_i(\mathbf{p})$. Examples: spheres, $SL(n)$ (dimension $n^2 - 1$, tangent space at $I$: traceless matrices), $O(n)$ (dimension $\frac{n(n-1)}{2}$, tangent space at $I$: skew-symmetric matrices).
    \item Lagrange multipliers say that $\nabla f$ is normal to the constraint manifold at a constrained extremum. The Riemannian gradient is the projection of $\nabla f$ onto the tangent space; Riemannian gradient descent alternates such steps with retractions.
    \item By Sard's theorem, almost all values of a smooth map are regular. A critical point is nondegenerate if its Hessian is invertible, and its index is the number of negative eigenvalues. Nondegenerate critical points are isolated, and by the Morse lemma $f$ is exactly a sum of $\pm y_i^2$ near them. Almost every linear perturbation of a function is a Morse function.
    \item Sublevel sets change their topology only at critical values: a critical point of index $k$ attaches a $k$-cell. On a compact manifold, $\chi(M) = \sum_k (-1)^k c_k$.
\end{itemize}

\section*{Exercises}
\addcontentsline{toc}{section}{Exercises}
\begin{enumerate}[leftmargin=*]
    \item Which of the following subsets of $\mathbb{R}^2$ are open, closed, both or neither? Find their interior, closure and boundary: (a) $\{x > 0, y > 0\}$; (b) $\{x^2 + y^2 \leq 1\}$; (c) $\{(x, 0) : 0 \leq x < 1\}$; (d) $\mathbb{Q}^2$, the points with rational coordinates.
    \item Show that the union of infinitely many closed sets need not be closed. Show that a set $U \subseteq X$ is open if and only if it is a union of open balls.
    \item \begin{enumerate}[label=(\alph*)]
        \item Show that the set of $m \times n$ matrices of rank at most $r$ is closed in $\mathbb{R}^{m \times n}$. (Hint: the rank is at most $r$ if and only if all $(r + 1) \times (r + 1)$ minors are $0$.)
        \item Show that every square matrix $A$ is a limit of invertible matrices. (Hint: $\det(A + tI)$ is a polynomial in $t$.) So $GL(n, \mathbb{R})$ is open and dense in $\mathbb{R}^{n \times n}$.
    \end{enumerate}
    \item Which of the following sets are compact? (a) The closed unit ball of the $\ell_1$ norm in $\mathbb{R}^n$; (b) $\{(x, y) : xy = 1\}$; (c) $SO(n)$; (d) $GL(n, \mathbb{R})$; (e) the ellipsoid $\{\mathbf{x} : \mathbf{x}^\top A\mathbf{x} \leq 1\}$ for a symmetric positive definite $A$; (f) the same set for $A = \diag(1, -1)$.
    \item Show that the regularized logistic regression loss $\frac{1}{N}\sum_i \log(1 + e^{-y_i\mathbf{w}^\top\mathbf{x}_i}) + \lambda\norm{\mathbf{w}}^2$ with $\lambda > 0$ has a global minimum, and that it is unique. (Use Proposition~\ref{A-prop:coercive} and Lecture~3.)
    \item Let $A \in \mathbb{R}^{n \times n}$. Show that $N(\mathbf{x}) = \norm{A\mathbf{x}}$ is a norm if and only if $A$ is invertible, and that the best constants in $c\norm{\mathbf{x}} \leq N(\mathbf{x}) \leq C\norm{\mathbf{x}}$ are the smallest and the largest singular value of $A$ (Lecture~2).
    \item \begin{enumerate}[label=(\alph*)]
        \item Show that the unit circle $S^1$ is path-connected, and that the hyperbola $\{x^2 - y^2 = 1\}$ is not connected.
        \item Show that $SO(2)$ is path-connected, and that $O(2)$ is not connected.
        \item Show that no continuous injective map from $S^1$ to $\mathbb{R}$ exists. (Hint: intermediate value theorem.)
    \end{enumerate}
    \item Use the iteration $x_{k+1} = \cos x_k$ with $x_0 = 1$ to compute the solution of $\cos x = x$ to six decimal places. How many steps does the a priori bound of Theorem~\ref{A-thm:banach} require, and how many do you actually need?
    \item Show that $f(x) = x + \frac{1}{2}\tanh(x)$ is a homeomorphism of $\mathbb{R}$, and compute $f^{-1}(1)$ with the iteration of Corollary~\ref{A-cor:residual}, starting at $x_0 = 0$. Is $f$ a diffeomorphism?
    \item At which points of the folium $x^3 + y^3 = 3xy$ can the curve not be written locally as a graph $y = g(x)$? Compute $g'(x)$ at the point $(\frac{3}{2}, \frac{3}{2})$.
    \item Let $F(x, y, z) = x^2 + y^2 - z^2$. Find all critical points and critical values of $F$. For which $c$ is $F^{-1}(c)$ a submanifold? Describe the level sets for $c > 0$, $c = 0$ and $c < 0$; which of them are connected?
    \item \begin{enumerate}[label=(\alph*)]
        \item Find the tangent space of the sphere $x^2 + y^2 + z^2 = 9$ at $(1, 2, 2)$ as $\Null DF$, and the tangent plane $(1, 2, 2) + T_{(1, 2, 2)}M$ as an equation.
        \item Show that $T_Q O(n) = \{QK : K^\top = -K\}$ and $T_I SL(n, \mathbb{R}) = \{H : \tr H = 0\}$ have the dimensions predicted by Example~\ref{A-ex:submanifolds}.
    \end{enumerate}
    \item (Orthogonal Procrustes problem.) Given $A \in \mathbb{R}^{n \times n}$ with SVD $A = U\Sigma V^\top$, show that $\tr(Q^\top A)$ is maximal over $Q \in O(n)$ for $Q = UV^\top$. (Hint: $\tr(Q^\top A) = \tr(\Sigma Z)$ with the orthogonal matrix $Z = V^\top Q^\top U$, whose entries are at most $1$ in absolute value.) This is how one aligns two point clouds by a rotation or reflection.
    \item \begin{enumerate}[label=(\alph*)]
        \item Find all critical points of the monkey saddle $f(x, y) = x^3 - 3xy^2$. Show that the origin is a degenerate critical point. In polar coordinates, $f = r^3\cos 3\varphi$; describe the regions where $f > 0$ and $f < 0$.
        \item For which $a \in \mathbb{R}$ is $x^3 - ax$ a Morse function? Find its critical points and their indices.
        \item Find the critical points and their indices of $f(x, y) = x^4 - 2x^2 + y^2 + \varepsilon x$ for a small $\varepsilon > 0$, and describe how the sublevel sets change as $c$ increases.
    \end{enumerate}
    \item Let $A = \diag(3, 2, 1)$ and $f(\mathbf{x}) = \mathbf{x}^\top A\mathbf{x}$ on the sphere $S^2$. Find all critical points of $f$ and their indices, and check the Morse relation $\chi(S^2) = 2$. Repeat for $A = \diag(4, 3, 2, 1)$ on $S^3$, whose Euler characteristic is $0$.
\end{enumerate}

\section*{Programming Exercises}
\addcontentsline{toc}{section}{Programming Exercises}

% Inline code cannot be hyphenated, so allow a little extra stretch in this section.
\setlength{\emergencystretch}{3em}

As in Lecture~3, use NumPy (and Matplotlib for plots), and start each file with
\begin{lstlisting}
import numpy as np
import matplotlib.pyplot as plt
rng = np.random.default_rng(0)
\end{lstlisting}

\begin{enumerate}[label=\textbf{P\arabic*.}, leftmargin=*]
    \item \textbf{Inverting a residual layer.} Let \py{W = rng.normal(size=(5, 5))} and $\mathbf{g}(\mathbf{x}) = c\tanh(W\mathbf{x})$ with $c = 0.9/\norm{W}_2$.
    \begin{itemize}
        \item Show that $\mathbf{g}$ is Lipschitz continuous with constant $0.9$ (use $|\tanh'| \leq 1$ and Lemma~\ref{A-lem:mean-value}).
        \item For a random $\mathbf{x}$, compute $\mathbf{y} = \mathbf{x} + \mathbf{g}(\mathbf{x})$, and recover $\mathbf{x}$ with the iteration of Corollary~\ref{A-cor:residual}. Plot the error against the number of iterations on a logarithmic scale, together with the bound of Theorem~\ref{A-thm:banach}.
        \item Repeat with $c = 1.5/\norm{W}_2$. Does the iteration still converge? Is $\mathbf{x} \mapsto \mathbf{x} + \mathbf{g}(\mathbf{x})$ still injective? (Try to find two points with the same image.)
    \end{itemize}

    \item \textbf{Riemannian gradient ascent on the sphere.} For a random symmetric $A = B + B^\top$ with \py{B = rng.normal(size=(6, 6))}, maximize $f(\mathbf{x}) = \mathbf{x}^\top A\mathbf{x}$ on $S^5$ with the Riemannian gradient and the normalization retraction (Section~\ref{A-sec:lagrange}).
    \begin{itemize}
        \item Start from $100$ random points, and check that you always reach $\pm\mathbf{q}_1$ (compare with \py{np.linalg.eigh}). Plot $f(\mathbf{p}_k)$ for a few runs.
        \item Start exactly at an eigenvector $\mathbf{q}_i$ with $i > 1$. What happens? Now add a perturbation of size $10^{-8}$. How many steps does it take to escape from the saddle point, and how does this depend on $\lambda_1 - \lambda_i$?
    \end{itemize}

    \item \textbf{Orthogonal Procrustes.} Generate a point cloud \py{X = rng.normal(size=(50, 3))}, a random orthogonal matrix $Q_0$ (from \py{np.linalg.qr}), and $Y = XQ_0^\top + 0.01\cdot$noise. Recover $Q_0$ by maximizing $\tr(Q^\top Y^\top X)$ over $O(3)$ with the SVD formula of Exercise~13. Then solve the same problem with Riemannian gradient ascent on $O(3)$: project $G = \nabla_Q$ onto $T_Q O(3)$ by $Q\,\frac{1}{2}(Q^\top G - G^\top Q)$, take a step, and retract with the $Q$ factor of a QR decomposition (fixing the signs so that the diagonal of $R$ is positive). Compare the results.

    \item \textbf{Counting components of sublevel sets.} For $f(x, y) = x^4 - 2x^2 + y^2 + 0.3x$, evaluate $f$ on a fine grid over $[-2, 2]^2$, and for $50$ values of $c$, count the connected components of the grid points with $f \leq c$ (write a breadth-first search over neighboring grid points). Plot the number of components against $c$, and compare the values of $c$ where it changes with the critical values of $f$, which you can find with Newton's method applied to $\nabla f = \mathbf{0}$.
\end{enumerate}
\endgroup
